\section{Countability and the First Shock of Infinity}
\chapterline

For centuries, mathematics treated infinity as a single concept: an endless process, written with the symbol $\infty$.

Then, in the 1870s, Georg Cantor asked a basic question:
\begin{center}
\emph{Can one infinity be strictly larger than another?}
\end{center}

To answer that question, Cantor realized we cannot count an infinite set to its end. Instead, we must return to functions: \textbf{bijective pairing}.

\begin{historybox}
In 1874, Georg Cantor published a paper proving that while the algebraic numbers (roots of polynomials with integer coefficients) can be placed in one-to-one correspondence with the natural numbers $\N$, the real numbers $\R$ cannot. This yielded an unexpected consequence: transcendental numbers like $\pi$ and $e$ were not rare exceptions; they formed an uncountable majority of the continuum.

Cantor's work initially met strong resistance. Leopold Kronecker, a strict finitist who remarked ``God made the integers, all else is human work,'' viewed transfinite numbers skeptically. Henri Poincar\'e famously expressed reservations about transfinite set theory.

In 1891, Cantor introduced his diagonal argument, proving that every set has a strictly larger power set ($|S| < |\pow(S)|$). David Hilbert later championed Cantor's work, declaring in 1926: ``No one shall expel us from the paradise that Cantor has created for us.''
\end{historybox}

\subsection{Measuring Size Without Counting}

Imagine a theater filled with spectators and seats. If every person is sitting in exactly one seat, and no seats are empty, you know with absolute certainty that the number of spectators equals the number of seats. You do not need to count whether there are $50$ people or $50{,}000$; the existence of a bijection proves their quantities match.

Cantor applied this exact principle to infinite sets:

\begin{definition}[Equinumerosity and Countability]
Two sets $A$ and $B$ have the \textbf{same cardinality} (written $|A| = |B|$) if there exists a bijection $f: A \to B$. We say $A$ and $B$ are \textbf{equinumerous}.

A set $S$ is \textbf{countably infinite} if there exists a bijection $f: \N \to S$. A set is \textbf{countable} if it is either finite or countably infinite. A set that is not countable is called \textbf{uncountable}.
\end{definition}

Notice the word \emph{cardinality}: it is simply the mathematician's formal name for the size of a set when measured by pairing. In plain terms, a set is countable if you can line up its elements in an infinite queue:
\[
s_1, \; s_2, \; s_3, \; s_4, \; \dots
\]
such that every single member of the set eventually appears at some finite position in the list.

\begin{recoverybox}
\textbf{Mental Check:} If someone claims a set $S$ is countably infinite, what concrete mathematical object are they obligated to produce?
\par\vspace{0.3em}
\emph{Answer:} An infinite numbered list $s_1, s_2, s_3, \dots$ that contains every single element of $S$ at some finite position $n \in \N$, with no omissions and no duplicates.
\end{recoverybox}

\subsection{Galileo's Paradox: The Evens Match the Naturals}

In 1638, in his final work \emph{Discorsi e dimostrazioni matematiche intorno a due nuove scienze}, Galileo Galilei pointed out a deep riddle about infinite collections. Consider the set of all positive even integers:
\[
2\N = \{2, 4, 6, 8, 10, \dots\}.
\]
Look at them laid out on the number line. Between every two even numbers sits an odd number that was skipped. Our spatial intuition insists there are only half as many even numbers as natural numbers $\N = \{1, 2, 3, \dots\}$.

Yet if we test their cardinality by pairing them up, something startling happens. Consider the function:
\[
f: \N \to 2\N, \quad f(n) = 2n.
\]
\begin{itemize}
\item If $2a = 2b$, then dividing by $2$ gives $a = b$ (injective: no two different counting numbers share the same even partner).
\item Every even integer $2k \in 2\N$ has an input $k \in \N$ where $f(k) = 2k$ (surjective: every even number is claimed).
\end{itemize}
The function $f$ is a bijection. Despite omitting all odd integers, the set of even numbers has the same cardinality as the set of all natural numbers:
\[
|2\N| = |\N|.
\]

\begin{center}
\begin{tikzpicture}[scale=0.90, >=Latex]
  \begin{scope}[yshift=2.2cm]
    \node[anchor=west, font=\footnotesize\bfseries] at (-0.3, 1.3) {1. The Spatial View: Evens Appear to be Only Half of the Integers};
    \draw[rounded corners=4pt, fill=gray!3, draw=gray!25] (-0.5, -0.6) rectangle (12.2, 1.1);
    
    \draw[thick, gray!60, ->] (0.2, 0.0) -- (11.5, 0.0);
    \foreach \x in {1.0, 2.2, 3.4, 4.6, 5.8, 7.0, 8.2, 9.4} {
      \draw[thick, gray!60] (\x, -0.08) -- (\x, 0.08);
    }
    
    \node[circle, draw=gray!40, fill=gray!10, inner sep=2pt, font=\scriptsize, text=gray!60] at (1.0, 0.42) {$1$};
    \node[circle, draw=targetgreen, fill=targetgreen!20, inner sep=2pt, font=\scriptsize\bfseries] at (2.2, 0.42) {$2$};
    \node[circle, draw=gray!40, fill=gray!10, inner sep=2pt, font=\scriptsize, text=gray!60] at (3.4, 0.42) {$3$};
    \node[circle, draw=targetgreen, fill=targetgreen!20, inner sep=2pt, font=\scriptsize\bfseries] at (4.6, 0.42) {$4$};
    \node[circle, draw=gray!40, fill=gray!10, inner sep=2pt, font=\scriptsize, text=gray!60] at (5.8, 0.42) {$5$};
    \node[circle, draw=targetgreen, fill=targetgreen!20, inner sep=2pt, font=\scriptsize\bfseries] at (7.0, 0.42) {$6$};
    \node[circle, draw=gray!40, fill=gray!10, inner sep=2pt, font=\scriptsize, text=gray!60] at (8.2, 0.42) {$7$};
    \node[circle, draw=targetgreen, fill=targetgreen!20, inner sep=2pt, font=\scriptsize\bfseries] at (9.4, 0.42) {$8$};
    \node[font=\scriptsize, text=gray!60] at (10.5, 0.42) {$\dots$};
    
    \node[anchor=north, font=\tiny\itshape, text=trapred!90!black] at (1.0, -0.15) {gap};
    \node[anchor=north, font=\tiny\itshape, text=trapred!90!black] at (3.4, -0.15) {gap};
    \node[anchor=north, font=\tiny\itshape, text=trapred!90!black] at (5.8, -0.15) {gap};
    \node[anchor=north, font=\tiny\itshape, text=trapred!90!black] at (8.2, -0.15) {gap};
  \end{scope}

  \begin{scope}[yshift=-1.2cm]
    \node[anchor=west, font=\footnotesize\bfseries] at (-0.3, 2.3) {2. The Bijection View: Every Number Finds a Partner ($|2\N| = |\N|$)};
    \draw[rounded corners=4pt, fill=gray!3, draw=gray!25] (-0.5, -0.7) rectangle (12.2, 2.1);
    
    \node[anchor=east, font=\scriptsize\bfseries, text=sourceblue!90!black] at (0.7, 1.4) {$\N:$};
    \node[circle, draw=sourceblue, fill=sourceblue!15, inner sep=2pt, font=\scriptsize\bfseries] (n1) at (1.4, 1.4) {$1$};
    \node[circle, draw=sourceblue, fill=sourceblue!15, inner sep=2pt, font=\scriptsize\bfseries] (n2) at (2.8, 1.4) {$2$};
    \node[circle, draw=sourceblue, fill=sourceblue!15, inner sep=2pt, font=\scriptsize\bfseries] (n3) at (4.2, 1.4) {$3$};
    \node[circle, draw=sourceblue, fill=sourceblue!15, inner sep=2pt, font=\scriptsize\bfseries] (n4) at (5.6, 1.4) {$4$};
    \node[circle, draw=sourceblue, fill=sourceblue!15, inner sep=2pt, font=\scriptsize\bfseries] (n5) at (7.0, 1.4) {$5$};
    \node[circle, draw=sourceblue, fill=sourceblue!15, inner sep=2pt, font=\scriptsize\bfseries] (n6) at (8.4, 1.4) {$6$};
    \node[font=\scriptsize, text=gray!60] (ndots) at (9.6, 1.4) {$\dots$};
    \node[font=\scriptsize, text=sourceblue!90!black] (nn) at (10.8, 1.4) {$n$};

    \node[anchor=east, font=\scriptsize\bfseries, text=targetgreen!90!black] at (0.7, 0.0) {$2\N:$};
    \node[circle, draw=targetgreen, fill=targetgreen!20, inner sep=2pt, font=\scriptsize\bfseries] (e1) at (1.4, 0.0) {$2$};
    \node[circle, draw=targetgreen, fill=targetgreen!20, inner sep=2pt, font=\scriptsize\bfseries] (e2) at (2.8, 0.0) {$4$};
    \node[circle, draw=targetgreen, fill=targetgreen!20, inner sep=2pt, font=\scriptsize\bfseries] (e3) at (4.2, 0.0) {$6$};
    \node[circle, draw=targetgreen, fill=targetgreen!20, inner sep=2pt, font=\scriptsize\bfseries] (e4) at (5.6, 0.0) {$8$};
    \node[circle, draw=targetgreen, fill=targetgreen!20, inner sep=2pt, font=\scriptsize\bfseries] (e5) at (7.0, 0.0) {$10$};
    \node[circle, draw=targetgreen, fill=targetgreen!20, inner sep=2pt, font=\scriptsize\bfseries] (e6) at (8.4, 0.0) {$12$};
    \node[font=\scriptsize, text=gray!60] (edots) at (9.6, 0.0) {$\dots$};
    \node[font=\scriptsize, text=targetgreen!90!black] (en) at (10.8, 0.0) {$2n$};

    \draw[<->, thick, quotientorange] (n1) -- (e1);
    \draw[<->, thick, quotientorange] (n2) -- (e2);
    \draw[<->, thick, quotientorange] (n3) -- (e3);
    \draw[<->, thick, quotientorange] (n4) -- (e4);
    \draw[<->, thick, quotientorange] (n5) -- (e5);
    \draw[<->, thick, quotientorange] (n6) -- (e6);
    \draw[<->, thick, quotientorange] (nn) -- node[right, font=\tiny\bfseries, text=quotientorange!90!black] {$f(n) = 2n$} (en);

    \node[anchor=north, font=\tiny\itshape, text=gray!70] at (5.6, -0.3) {Every spectator has a seat; no seats are empty. The sizes are identical.};
  \end{scope}
\end{tikzpicture}
\end{center}

For infinite sets, Euclid's classical common notion ``the whole is greater than the part'' ceases to hold. In fact, Richard Dedekind took this as the very defining property of the infinite: \emph{a set $S$ is infinite if and only if it can be placed in one-to-one correspondence with a proper subset of itself.} For any finite set, the Pigeonhole Principle guarantees that removing an element strictly decreases the size; for an infinite set, you can remove infinitely many elements and still have just as many left over.

\begin{recoverybox}
\textbf{Mental Check:} Can you construct an explicit bijection between the natural numbers $\N$ and the positive \emph{odd} integers $\{1, 3, 5, 7, \dots\}$?
\par\vspace{0.3em}
\emph{Answer:} Yes! The function $g(n) = 2n - 1$ sends $1 \mapsto 1$, $2 \mapsto 3$, $3 \mapsto 5$, pairing every natural number with a unique odd number.
\end{recoverybox}

\subsection{The Integers are Countable}

What about the set of all integers $\Z = \{\dots, -3, -2, -1, 0, 1, 2, 3, \dots\}$? 

Unlike $\N$, which begins cleanly at $1$, the integers stretch endlessly in two opposite directions. If you tried to count them by walking out along the positives $0, 1, 2, 3, \dots$, you would walk forever and never reach $-1$. 

The solution is to weave the two directions together, alternating between positive and negative values:

\begin{center}
\begin{tikzpicture}[scale=0.92, >=Latex]
  \draw[rounded corners=4pt, fill=gray!3, draw=gray!25] (-5.8, -2.8) rectangle (5.8, 2.6);
  \node[anchor=south west, font=\footnotesize\bfseries] at (-5.6, 2.2) {The Outward Spiral: Enumerating the Two-Way Infinite Integers $\Z$};

  % Integer line segments between nodes (no line passing through any text)
  \draw[thick, gray!60, <-] (-5.2, 0.0) -- (-4.3, 0.0);
  \draw[thick, gray!60] (-3.5, 0.0) -- (-3.0, 0.0);
  \draw[thick, gray!60] (-2.2, 0.0) -- (-1.7, 0.0);
  \draw[thick, gray!60] (-0.9, 0.0) -- (-0.4, 0.0);
  \draw[thick, gray!60] (0.4, 0.0) -- (0.9, 0.0);
  \draw[thick, gray!60] (1.7, 0.0) -- (2.2, 0.0);
  \draw[thick, gray!60] (3.0, 0.0) -- (3.5, 0.0);
  \draw[thick, gray!60, ->] (4.3, 0.0) -- (5.2, 0.0);
  
  \node[circle, draw=targetgreen, fill=targetgreen!20, inner sep=2pt, font=\scriptsize\bfseries] (z0) at (0, 0) {$0$};
  \node[circle, draw=sourceblue, fill=sourceblue!15, inner sep=2pt, font=\scriptsize\bfseries] (zp1) at (1.3, 0) {$1$};
  \node[circle, draw=quotientorange, fill=quotientorange!20, inner sep=2pt, font=\scriptsize\bfseries] (zm1) at (-1.3, 0) {$-1$};
  \node[circle, draw=sourceblue, fill=sourceblue!15, inner sep=2pt, font=\scriptsize\bfseries] (zp2) at (2.6, 0) {$2$};
  \node[circle, draw=quotientorange, fill=quotientorange!20, inner sep=2pt, font=\scriptsize\bfseries] (zm2) at (-2.6, 0) {$-2$};
  \node[circle, draw=sourceblue, fill=sourceblue!15, inner sep=2pt, font=\scriptsize\bfseries] (zp3) at (3.9, 0) {$3$};
  \node[circle, draw=quotientorange, fill=quotientorange!20, inner sep=2pt, font=\scriptsize\bfseries] (zm3) at (-3.9, 0) {$-3$};
  
  % Step 1: at 0
  \node[anchor=south, font=\tiny\bfseries, text=targetgreen!80!black] at (0, 0.35) {step $1$};
  
  % Step 2: 0 -> +1 (arc above)
  \draw[->, thick, structurepurple] (0.2, 0.25) to[bend left=55] node[pos=0.75, above right=-1pt, font=\tiny\bfseries, text=structurepurple] {$2$} (1.1, 0.25);
  
  % Step 3: +1 -> -1 (arc below)
  \draw[->, thick, structurepurple] (1.1, -0.25) to[bend left=42] node[pos=0.8, below left=-1pt, font=\tiny\bfseries, text=structurepurple] {$3$} (-1.1, -0.25);
  
  % Step 4: -1 -> +2 (arc above)
  \draw[->, thick, structurepurple] (-1.1, 0.25) to[bend left=45] node[pos=0.8, above right=-1pt, font=\tiny\bfseries, text=structurepurple] {$4$} (2.4, 0.25);
  
  % Step 5: +2 -> -2 (arc below)
  \draw[->, thick, structurepurple] (2.4, -0.25) to[bend left=40] node[pos=0.8, below left=-1pt, font=\tiny\bfseries, text=structurepurple] {$5$} (-2.4, -0.25);
  
  % Step 6: -2 -> +3 (arc above)
  \draw[->, thick, structurepurple] (-2.4, 0.25) to[bend left=38] node[pos=0.85, above right=-1pt, font=\tiny\bfseries, text=structurepurple] {$6$} (3.7, 0.25);
  
  % Step 7: +3 -> -3 (arc below)
  \draw[->, thick, structurepurple] (3.7, -0.25) to[bend left=36] node[pos=0.85, below left=-1pt, font=\tiny\bfseries, text=structurepurple] {$7$} (-3.7, -0.25);

  % Bottom queue representation
  \node[anchor=north, font=\scriptsize, align=center] at (0, -1.75) {
    \textbf{Resulting Queue in $\N$:}\\[0.25em]
    $f(1)=0, \; f(2)=1, \; f(3)=-1, \; f(4)=2, \; f(5)=-2, \; f(6)=3, \dots$
  };
\end{tikzpicture}\\[0.5em]
{\small\textbf{Figure 19.2:}\label{19:fig-integers-spiral} \emph{The outward spiral bijection between $\N$ and $\Z$. Walking in a single direction would abandon all negative numbers. By spiraling outward from $0$ in alternating steps ($+1, -1, +2, -2, \dots$), every integer is reached at a predictable, finite index, proving $|\Z| = |\N|$.}}
\end{center}
The explicit bijection $f: \N \to \Z$ is:
\[
f(n) = \begin{cases} \frac{n}{2} & \text{if } n \text{ is even}, \\ -\frac{n-1}{2} & \text{if } n \text{ is odd}. \end{cases}
\]
To verify that this really is a bijection, we can write down its exact inverse $f^{-1}: \Z \to \N$. If you hand me any integer $z$, I can instantly tell you its position $n$ in line:
\[
f^{-1}(z) = \begin{cases} 2z & \text{if } z > 0, \\ 1 - 2z & \text{if } z \le 0. \end{cases}
\]
For example, the integer $-3$ appears at index $1 - 2(-3) = 7$, exactly matching our table. Every integer appears at a predictable, finite index. Thus $|\Z| = |\N|$: the two-way infinite line has no more elements than the one-way ray!

\subsection{The Rationals are Countable: Cantor's Snake}

The rational numbers $\Q$ seem much larger than the integers. Between any two integers lies an infinite collection of fractions. Even more strikingly, between any two fractions $a < b$, you can always find another fraction sitting between them---for instance, their midpoint $(a+b)/2$. And between that midpoint and either end, there are still more fractions, ad infinitum. 

Unlike the integers, where $4$ is followed immediately by $5$ with nothing in between, the rationals have no ``next-door neighbor.'' If you stand at $1/2$, what is the next rational number? There is no such thing! Mathematicians say the rational numbers are \textbf{dense} on the real line. 

If you try to list the rationals in increasing numerical order, there are infinitely many fractions between $0$ and $1$, so you would never even reach $2$.

Yet density is a topological property describing how numbers are arranged in space, not a measure of cardinality: the rationals are countable.

To see why, Cantor realized we must abandon the number line entirely. Instead of ordering fractions by their numerical value, arrange all pairs of positive integers $(p, q)$ in an infinite two-dimensional grid, where $p$ is the row (the numerator) and $q$ is the column (the denominator). We then weave through the grid along finite diagonals:

\begin{center}
\begin{tikzpicture}[scale=0.92, >=Latex]
  \node[anchor=west, font=\footnotesize\bfseries] at (-0.3, 3.8) {Cantor's Diagonal Snake: Enumerating the Rationals $\N \times \N$};
  \draw[rounded corners=4pt, fill=gray!3, draw=gray!25] (-0.5, -0.6) rectangle (12.2, 3.5);
  
  % Column labels
  \node[font=\scriptsize\bfseries, text=gray!70] at (1.4, 3.1) {$q=1$};
  \node[font=\scriptsize\bfseries, text=gray!70] at (2.8, 3.1) {$q=2$};
  \node[font=\scriptsize\bfseries, text=gray!70] at (4.2, 3.1) {$q=3$};
  \node[font=\scriptsize\bfseries, text=gray!70] at (5.6, 3.1) {$q=4$};
  
  % Row labels
  \node[font=\scriptsize\bfseries, text=gray!70] at (0.3, 2.4) {$p=1$};
  \node[font=\scriptsize\bfseries, text=gray!70] at (0.3, 1.6) {$p=2$};
  \node[font=\scriptsize\bfseries, text=gray!70] at (0.3, 0.8) {$p=3$};
  \node[font=\scriptsize\bfseries, text=gray!70] at (0.3, 0.0) {$p=4$};
  
  % Nodes
  \node[circle, draw=sourceblue, fill=sourceblue!15, inner sep=2pt, font=\scriptsize] (f11) at (1.4, 2.4) {$\frac{1}{1}$};
  \node[circle, draw=sourceblue, fill=sourceblue!15, inner sep=2pt, font=\scriptsize] (f12) at (2.8, 2.4) {$\frac{1}{2}$};
  \node[circle, draw=sourceblue, fill=sourceblue!15, inner sep=2pt, font=\scriptsize] (f13) at (4.2, 2.4) {$\frac{1}{3}$};
  \node[circle, draw=sourceblue, fill=sourceblue!15, inner sep=2pt, font=\scriptsize] (f14) at (5.6, 2.4) {$\frac{1}{4}$};

  \node[circle, draw=sourceblue, fill=sourceblue!15, inner sep=2pt, font=\scriptsize] (f21) at (1.4, 1.6) {$\frac{2}{1}$};
  \node[circle, draw=gray!50, fill=gray!15, inner sep=2pt, font=\scriptsize, text=gray!60] (f22) at (2.8, 1.6) {$\frac{2}{2}$};
  \node[circle, draw=sourceblue, fill=sourceblue!15, inner sep=2pt, font=\scriptsize] (f23) at (4.2, 1.6) {$\frac{2}{3}$};

  \node[circle, draw=sourceblue, fill=sourceblue!15, inner sep=2pt, font=\scriptsize] (f31) at (1.4, 0.8) {$\frac{3}{1}$};
  \node[circle, draw=sourceblue, fill=sourceblue!15, inner sep=2pt, font=\scriptsize] (f32) at (2.8, 0.8) {$\frac{3}{2}$};

  \node[circle, draw=sourceblue, fill=sourceblue!15, inner sep=2pt, font=\scriptsize] (f41) at (1.4, 0.0) {$\frac{4}{1}$};
  
  % Order Badges
  \node[font=\tiny\bfseries, text=quotientorange!90!black, anchor=south east] at (f11.north west) {1};
  \node[font=\tiny\bfseries, text=quotientorange!90!black, anchor=south] at (f12.north) {2};
  \node[font=\tiny\bfseries, text=quotientorange!90!black, anchor=east] at (f21.west) {3};
  \node[font=\tiny\bfseries, text=quotientorange!90!black, anchor=east] at (f31.west) {4};
  \node[font=\tiny\bfseries, text=trapred!80!black, anchor=west] at (f22.east) {skip};
  \node[font=\tiny\bfseries, text=quotientorange!90!black, anchor=south] at (f13.north) {5};
  \node[font=\tiny\bfseries, text=quotientorange!90!black, anchor=south] at (f14.north) {6};
  \node[font=\tiny\bfseries, text=quotientorange!90!black, anchor=west] at (f23.east) {7};
  \node[font=\tiny\bfseries, text=quotientorange!90!black, anchor=west] at (f32.east) {8};
  \node[font=\tiny\bfseries, text=quotientorange!90!black, anchor=north] at (f41.south) {9};
  
  % Snake traversal arrows
  \draw[->, thick, quotientorange] (f11) -- (f12);
  \draw[->, thick, quotientorange] (f12) -- (f21);
  \draw[->, thick, quotientorange] (f21) -- (f31);
  \draw[->, thick, quotientorange] (f31) -- (f22);
  \draw[->, thick, quotientorange] (f22) -- (f13);
  \draw[->, thick, quotientorange] (f13) -- (f14);
  \draw[->, thick, quotientorange] (f14) -- (f23);
  \draw[->, thick, quotientorange] (f23) -- (f32);
  \draw[->, thick, quotientorange] (f32) -- (f41);
  
  % Side panel description
  \node[anchor=west, font=\scriptsize, align=left] at (6.8, 1.45) {
    \textbf{Traversal by Diagonals:}\\[0.15em]
    $\bullet$ Diagonal $p+q = 2$: $\{1/1\}$ (1 entry)\\[0.1em]
    $\bullet$ Diagonal $p+q = 3$: $\{1/2, 2/1\}$ (2 entries)\\[0.1em]
    $\bullet$ Diagonal $p+q = 4$: $\{3/1, 2/2, 1/3\}$ (3 entries)\\[0.1em]
    $\bullet$ Diagonal $p+q = 5$: $\{1/4, 2/3, 3/2, 4/1\}$ (4 entries)\\[0.35em]
    \textbf{Why it never gets trapped:}\\[0.15em]
    Every fraction $p/q$ has a finite sum $p+q = k$.\\
    Each diagonal is finite, so step $k$ is\\
    reached after finitely many rungs!\\
    Skipping unreduced fractions ($2/2 = 1/1$)\\
    yields an exact bijection $\N \to \Q^+$.
  };
\end{tikzpicture}\\[0.5em]
{\small\textbf{Figure 19.3:}\label{19:fig-cantor-snake} \emph{Cantor's diagonal snake enumerating the positive rationals $\Q^+$. By weaving along finite diagonals $p+q = k$, every fraction is reached after finitely many steps, proving that the dense rationals are countable ($|\Q| = |\N|$).}}
\end{center}

Because each diagonal $p + q = k$ contains only $k - 1$ entries (a finite number), the snake reaches any chosen fraction $p/q$ after a finite number of steps. For instance, the fraction $355/113$ sits on diagonal $355 + 113 = 468$; because all previous $466$ diagonals together contain only finitely many fractions, the snake arrives at $355/113$ after finitely many steps.

By omitting duplicate fractions that are not in lowest terms (such as $2/2 = 1/1$, $4/2 = 2/1$, or $2/4 = 1/2$), each positive rational number is visited exactly once. This establishes an explicit bijection between $\N$ and the positive rationals $\Q^+$.

Just as we did with the integers, we can now interleave the positive rationals with their negative counterparts and insert $0$ at the very front:
\[
0, \; \frac{1}{1}, \; -\frac{1}{1}, \; \frac{1}{2}, \; -\frac{1}{2}, \; \frac{2}{1}, \; -\frac{2}{1}, \; \frac{1}{3}, \; -\frac{1}{3}, \; \dots
\]
Every rational number in existence appears at a definite, finite position. Therefore:
\[
|\Q| = |\N|.
\]
Density does not make a set uncountable: even though fractions are densely packed along the number line, the set of all rational numbers can be arranged into a single infinite sequence.

\begin{recoverybox}
\textbf{Mental Check:} Why can't we count the grid by reading row 1 all the way across: $1/1, 1/2, 1/3, 1/4, \dots$?
\par\vspace{0.3em}
\emph{Answer:} Because row 1 is infinite! If you tried to finish row 1 before moving to row 2, you would walk forever along row 1 and never reach $2/1$. The diagonal snake works because every diagonal is finite.
\end{recoverybox}

\subsection{The Shock: The Real Numbers are Uncountable}

At this point, one might wonder: is \emph{every} infinite set countable? If both the integers and the dense set of rational numbers can be arranged into sequences, could it be that all infinite sets have the same cardinality?

Cantor showed that the answer is no: the real numbers cannot be listed in an infinite sequence.

\begin{theorem}[Cantor's Diagonal Theorem]
The open interval of real numbers $(0, 1)$ is uncountable.
\end{theorem}

\begin{center}
\begin{tikzpicture}[scale=0.92, >=Latex]
  \node[anchor=west, font=\footnotesize\bfseries] at (-0.3, 4.3) {Cantor's Diagonal Construction: A Real Missing from the Catalog};
  \draw[rounded corners=4pt, fill=gray!3, draw=gray!25] (-0.5, -0.8) rectangle (12.2, 4.0);
  
  % Rows of cataloged reals
  \node[font=\footnotesize\bfseries, text=sourceblue!90!black] at (0.8, 3.2) {$r_1 = 0.$};
  \node[draw=quotientorange, fill=quotientorange!25, rounded corners=1pt, inner sep=2pt, font=\footnotesize\bfseries] at (2.0, 3.2) {$7$};
  \node[font=\footnotesize] at (2.6, 3.2) {$1$};
  \node[font=\footnotesize] at (3.2, 3.2) {$4$};
  \node[font=\footnotesize] at (3.8, 3.2) {$2$};
  \node[font=\footnotesize] at (4.4, 3.2) {$8$};
  \node[font=\footnotesize, text=gray!50] at (5.2, 3.2) {$\dots$};

  \node[font=\footnotesize\bfseries, text=sourceblue!90!black] at (0.8, 2.5) {$r_2 = 0.$};
  \node[font=\footnotesize] at (2.0, 2.5) {$3$};
  \node[draw=quotientorange, fill=quotientorange!25, rounded corners=1pt, inner sep=2pt, font=\footnotesize\bfseries] at (2.6, 2.5) {$5$};
  \node[font=\footnotesize] at (3.2, 2.5) {$9$};
  \node[font=\footnotesize] at (3.8, 2.5) {$2$};
  \node[font=\footnotesize] at (4.4, 2.5) {$6$};
  \node[font=\footnotesize, text=gray!50] at (5.2, 2.5) {$\dots$};

  \node[font=\footnotesize\bfseries, text=sourceblue!90!black] at (0.8, 1.8) {$r_3 = 0.$};
  \node[font=\footnotesize] at (2.0, 1.8) {$0$};
  \node[font=\footnotesize] at (2.6, 1.8) {$8$};
  \node[draw=quotientorange, fill=quotientorange!25, rounded corners=1pt, inner sep=2pt, font=\footnotesize\bfseries] at (3.2, 1.8) {$3$};
  \node[font=\footnotesize] at (3.8, 1.8) {$3$};
  \node[font=\footnotesize] at (4.4, 1.8) {$3$};
  \node[font=\footnotesize, text=gray!50] at (5.2, 1.8) {$\dots$};

  \node[font=\footnotesize\bfseries, text=sourceblue!90!black] at (0.8, 1.1) {$r_4 = 0.$};
  \node[font=\footnotesize] at (2.0, 1.1) {$4$};
  \node[font=\footnotesize] at (2.6, 1.1) {$1$};
  \node[font=\footnotesize] at (3.2, 1.1) {$9$};
  \node[draw=quotientorange, fill=quotientorange!25, rounded corners=1pt, inner sep=2pt, font=\footnotesize\bfseries] at (3.8, 1.1) {$4$};
  \node[font=\footnotesize] at (4.4, 1.1) {$0$};
  \node[font=\footnotesize, text=gray!50] at (5.2, 1.1) {$\dots$};

  % Divider
  \draw[thick, gray!40] (0.2, 0.6) -- (5.6, 0.6);

  % Constructed rebel number d
  \node[font=\footnotesize\bfseries, text=targetgreen!90!black] at (0.8, 0.1) {$d \; = 0.$};
  \node[draw=targetgreen, fill=targetgreen!25, rounded corners=1pt, inner sep=2pt, font=\footnotesize\bfseries] at (2.0, 0.1) {$4$};
  \node[draw=targetgreen, fill=targetgreen!25, rounded corners=1pt, inner sep=2pt, font=\footnotesize\bfseries] at (2.6, 0.1) {$4$};
  \node[draw=targetgreen, fill=targetgreen!25, rounded corners=1pt, inner sep=2pt, font=\footnotesize\bfseries] at (3.2, 0.1) {$4$};
  \node[draw=targetgreen, fill=targetgreen!25, rounded corners=1pt, inner sep=2pt, font=\footnotesize\bfseries] at (3.8, 0.1) {$5$};
  \node[font=\footnotesize, text=gray!50] at (5.2, 0.1) {$\dots$};

  % Callout text
  \node[anchor=west, font=\scriptsize, align=left] at (6.0, 1.8) {
    $\bullet$ Diagonal digits $d_{nn}$ shown in orange.\\[0.5em]
    $\bullet$ Constructed digit $c_n$ (in green):\\
    \quad $c_n = 5$ if $d_{nn} = 4$; otherwise $c_n = 4$.\\[0.6em]
    $\bullet$ By construction, $d$ differs from $r_n$ at position $n$.\\[0.5em]
    $\bullet$ Hence $d \neq r_n$ for every $n$: $d$ is not on the list!
  };
\end{tikzpicture}\\[0.5em]
{\small\textbf{Figure 19.4:}\label{19:fig-cantor-diagonal} \emph{Cantor's diagonal construction. Given any purported list $r_1, r_2, \dots$ of all reals in $(0, 1)$, we construct a real number $d$ whose $n$-th digit differs from the diagonal digit $d_{nn}$. Since $d$ differs from $r_n$ at digit $n$ for every $n$, $d$ cannot appear on the list, proving $(0, 1)$ is uncountable.}}
\end{center}

\begin{proof}
We argue by contradiction. Suppose that $(0, 1)$ were countable. 
Then there exists a bijection from $\N$ to $(0, 1)$, meaning we can write down a complete, numbered catalog listing every real number in $(0, 1)$:
\[
r_1, \; r_2, \; r_3, \; r_4, \; \dots
\]
Let us express each number in its infinite decimal expansion:
\begin{align*}
r_1 &= 0.\mathbf{d_{11}} d_{12} d_{13} d_{14} \dots \\
r_2 &= 0.d_{21} \mathbf{d_{22}} d_{23} d_{24} \dots \\
r_3 &= 0.d_{31} d_{32} \mathbf{d_{33}} d_{34} \dots \\
r_4 &= 0.d_{41} d_{42} d_{43} \mathbf{d_{44}} \dots
\end{align*}
where $d_{ij} \in \{0, 1, \dots, 9\}$ is the $j$-th decimal digit of $r_i$.

We now construct a real number $d = 0.c_1 c_2 c_3 c_4 \dots$ that differs from every number on the list. For each integer $n \ge 1$, we inspect the diagonal digit $d_{nn}$ and choose $c_n$ according to the simple rule:
\[
c_n = \begin{cases} 5 & \text{if } d_{nn} = 4, \\ 4 & \text{if } d_{nn} \neq 4. \end{cases}
\]
Notice three properties of this constructed number $d$:
\begin{enumerate}[label=\arabic*.]
\item Because every digit $c_n$ is strictly either $4$ or $5$, $0 < d < 1$. Thus $d \in (0, 1)$.
\item The decimal expansion of $d$ contains only $4$s and $5$s, avoiding trailing infinite sequences of $0$s or $9$s. Its decimal representation is strictly unique.
\item Compare $d$ to any number $r_n$ on the list: their $n$-th decimal digits differ by construction ($c_n \neq d_{nn}$). Therefore:
\[
d \neq r_n \quad \text{for every integer } n \ge 1.
\]
\end{enumerate}

This is a direct contradiction! The list $r_1, r_2, \dots$ was assumed to contain \emph{every} real number in $(0, 1)$, yet the number $d$ lives in $(0, 1)$ and appears nowhere on the list.

Therefore, our initial assumption was false. The interval $(0, 1)$ cannot be listed in a sequence, and is uncountable.
\end{proof}

\begin{trapbox}
\textbf{The Trap of Dual Decimal Expansions}

Why did we strictly restrict our choices to the digits $4$ and $5$? Why not simply set $c_n = d_{nn} + 1 \pmod{10}$?

Because in base-10 arithmetic, some numbers possess \emph{two} distinct decimal representations:
\[
0.49999\dots = 0.50000\dots
\]
If our construction produced a number ending in infinite trailing $9$s (or infinite trailing $0$s), an argument could be raised that $d$ equals a listed number with an alternate representation (such as $0.4999\dots = 0.5000\dots$).

By choosing every digit $c_n$ strictly from $\{4, 5\}$, the number $d$ never ends in trailing $9$s or trailing $0$s. Its decimal expansion is unique, ensuring that $d \neq r_n$ holds unambiguously.
\end{trapbox}

Since $(0, 1) \subseteq \R$, the full set of real numbers $\R$ must also be uncountable. There are multiple, distinct tiers of infinity:
\[
|\N| = |\Z| = |\Q| < |\R|.
\]

\subsection{What Cardinality Unlocked, Its Limits, and Where We Go Next}

Cantor's cardinality theory led to fundamental developments across mathematics:
\begin{itemize}
\item \textbf{Non-constructive existence:} An \emph{algebraic number} is any root of a polynomial with integer coefficients (such as $\sqrt{2}$ or $-3/5$). Because polynomials of degree $n$ have only finitely many roots, and there are only countably many integer polynomials, the set of algebraic numbers is countable. But $\R$ is uncountable! This means the \emph{transcendental numbers} (numbers like $\pi$ and $e$ that cannot be solved by integer polynomials) must be uncountably infinite. Cantor proved that almost all real numbers are transcendental without constructing a single explicit example.
\item \textbf{The foundations of computing:} In 1936, Alan Turing used Cantor's diagonal argument to prove that the Halting Problem is uncomputable. Because computer programs (and Turing machines) are finite strings of text written in a finite alphabet, there are only countably many computer programs ($\aleph_0$). However, there are uncountably many mathematical functions from $\N$ to $\{0, 1\}$ ($2^{\aleph_0}$). Thus, the vast majority of problems can never be solved by any algorithm!
\item \textbf{The transfinite hierarchy:} Cantor proved that for \emph{any} set $S$, the power set $\pow(S)$ always has strictly greater cardinality than $S$ ($|S| < |\pow(S)|$). This established an infinite hierarchy of transfinite cardinalities: $\aleph_0 < 2^{\aleph_0} < 2^{2^{\aleph_0}} < \dots$, showing there is no maximal cardinality.
\end{itemize}

Yet cardinality does not capture geometric or topological structure. 

It counts elements by bijection, but does not distinguish dimension or distance. Consider this surprising fact, also discovered by Cantor:
\[
|(0, 1)| = |\R| = |\R^2| = |\R^{100}|.
\]
An interval, the real line, a two-dimensional plane, and a 100-dimensional space all have the same cardinality. When Cantor proved in 1877 that a line segment can be mapped bijectively to a unit square by interleaving digits, he famously wrote to Richard Dedekind:
\begin{center}
\emph{``Je le vois, mais je ne le crois pas!''} (``I see it, but I do not believe it!'')
\end{center}

\begin{center}
\begin{tikzpicture}[scale=0.92, >=Latex]
  \draw[rounded corners=4pt, fill=gray!3, draw=gray!25] (-5.5, -1.2) rectangle (5.5, 2.4);
  \node[anchor=south west, font=\footnotesize\bfseries] at (-5.3, 2.0) {Geometric Bijection: Projecting a Finite Interval to the Infinite Line $\R$};

  % Center of projection
  \coordinate (C) at (0, 1.5);
  \fill[targetgreen!80!black] (C) circle (2pt);
  \node[above=2pt, font=\scriptsize\bfseries, text=targetgreen!80!black] at (C) {Center of Projection $P$};

  % Bounded semicircular arc
  \draw[very thick, sourceblue] (-0.97, 1.24) arc[start angle=195, end angle=345, radius=1.0];
  \draw[fill=white, draw=sourceblue, thick] (-0.97, 1.24) circle (1.8pt);
  \draw[fill=white, draw=sourceblue, thick] (0.97, 1.24) circle (1.8pt);
  \node[anchor=west, font=\tiny\bfseries, text=sourceblue] at (1.1, 1.1) {interval $(-\pi/2, \pi/2)$};

  % Infinite real line
  \draw[thick, gray!70, <->] (-5.0, -0.4) -- (5.0, -0.4);
  \node[above=2pt, font=\scriptsize\bfseries, text=gray!80!black] at (4.7, -0.4) {$\R$};

  % Projection rays
  \draw[dashed, structurepurple!70!black] (C) -- (0, -0.4);
  \draw[dashed, structurepurple!70!black] (C) -- (1.58, -0.4);
  \draw[dashed, structurepurple!70!black] (C) -- (3.33, -0.4);
  \draw[dashed, structurepurple!70!black] (C) -- (-1.58, -0.4);
  \draw[dashed, structurepurple!70!black] (C) -- (-3.33, -0.4);

  % Points on arc
  \fill[sourceblue] (0, 0.5) circle (1.5pt);
  \fill[sourceblue] (0.64, 0.73) circle (1.5pt);
  \fill[sourceblue] (0.87, 1.00) circle (1.5pt);
  \fill[sourceblue] (-0.64, 0.73) circle (1.5pt);
  \fill[sourceblue] (-0.87, 1.00) circle (1.5pt);

  % Points on line
  \fill[targetgreen!80!black] (0, -0.4) circle (1.6pt) node[below=1pt, font=\tiny\bfseries] {$0$};
  \fill[targetgreen!80!black] (1.58, -0.4) circle (1.6pt) node[below=1pt, font=\tiny\bfseries] {$x_1$};
  \fill[targetgreen!80!black] (3.33, -0.4) circle (1.6pt) node[below=1pt, font=\tiny\bfseries] {$x_2$};
  \fill[targetgreen!80!black] (-1.58, -0.4) circle (1.6pt) node[below=1pt, font=\tiny\bfseries] {$-x_1$};
  \fill[targetgreen!80!black] (-3.33, -0.4) circle (1.6pt) node[below=1pt, font=\tiny\bfseries] {$-x_2$};

  \node[anchor=south, font=\tiny\itshape, text=structurepurple!90!black] at (0, -1.05) {Each ray through angle $\theta$ hits exactly one point $x = \tan\theta$, giving $|(-\pi/2, \pi/2)| = |\R|$.};
\end{tikzpicture}\\[0.5em]
{\small\textbf{Figure 19.5:}\label{19:fig-interval-to-real-line} \emph{Stereographic projection establishing equinumerosity $|(-\pi/2, \pi/2)| = |\R|$. Bending the open interval into an arc and casting rays from the center $P$ matches each angle $\theta$ with exactly one real number $x = \tan\theta$. The bounded interval has the exact same cardinality as the infinite real line.}}
\end{center}

Yet a plane is topologically and geometrically different from a line, and spaces of different dimensions possess distinct geometric properties. Cardinality does not detect dimension, distance, or continuity: it cannot determine whether a function is continuous or whether a sequence of points converges.

To do calculus and analysis, we cannot rely on cardinality alone. We must equip the continuum with a notion of distance and limits. In Chapter~20, we turn to \textbf{Real Analysis Preview: Epsilon as a Contract}, showing how Cauchy and Weierstrass formalized limits using quantifier definitions.

\subsection{Exercises}

\begin{enumerate}
\item \textbf{Mental Check: Finite Strings vs Infinite Sequences.}
Let $\Sigma = \{0, 1\}$.
\begin{enumerate}[label=(\alph*)]
\item Prove that the set of all \emph{finite} binary strings $\Sigma^*$ (strings of length $0, 1, 2, \dots$) is countably infinite. (Hint: List the empty string, then strings of length $1$, then strings of length $2$, and so on.)
\item Prove that the set of all \emph{infinite} binary sequences $\Sigma^\omega = \{(b_1, b_2, b_3, \dots) : b_i \in \{0, 1\}\}$ is uncountable by adapting Cantor's diagonal argument.
\end{enumerate}

\item \textbf{Rescaling the Continuum.}
\begin{enumerate}[label=(\alph*)]
\item For any two real numbers $a < b$, construct an explicit linear bijection $f: (0, 1) \to (a, b)$, proving that all open intervals have the same cardinality regardless of their length.
\item Using the tangent function, construct an explicit bijection $g: (-\pi/2, \pi/2) \to \R$, proving that a finite interval of length $\pi$ has the exact same cardinality as the entire infinite real line.
\end{enumerate}

\item \textbf{Countable Union of Countable Sets.}
Suppose $\{A_n\}_{n=1}^\infty$ is a countably infinite collection of sets, where each set $A_n = \{a_{n,1}, a_{n,2}, a_{n,3}, \dots\}$ is countably infinite. Prove that their union:
\[
A = \bigcup_{n=1}^\infty A_n
\]
is countably infinite. (Hint: Arrange the elements into a two-dimensional grid where row $n$ lists the elements of $A_n$, and adapt Cantor's diagonal snake traversal.)

\item \textbf{Cantor's General Theorem: Diagonalization in the Abstract.}
Prove that for \emph{any} set $S$ (finite or infinite), there is no surjective function from $S$ to its power set $\pow(S)$.
\begin{enumerate}[label=(\alph*)]
\item Suppose towards a contradiction that $f: S \to \pow(S)$ is surjective. Define the subset:
\[
D = \{x \in S : x \notin f(x)\}.
\]
\item Because $f$ is surjective and $D \in \pow(S)$, there must exist an element $d \in S$ such that $f(d) = D$.
\item Ask the critical question: Is $d \in D$? Show that $d \in D \implies d \notin D$, and $d \notin D \implies d \in D$, producing a contradiction.
\item Reflect: How is the definition of $D$ related to Cantor's diagonal number $d$, Russell's paradox, and the Halting Problem?
\end{enumerate}

\item \textbf{A Glimpse of Hilbert's Hotel on the Continuum.}
Prove that the half-open interval $(0, 1]$ has the exact same cardinality as the open interval $(0, 1)$ by constructing an explicit bijection $h: (0, 1] \to (0, 1)$.
\par\vspace{0.2em}
\emph{Hint:} Isolate a countable sequence in $(0, 1]$, say $S = \{1, 1/2, 1/4, \dots, 1/2^k, \dots\}$. Shift each element of $S$ to the next power of $2$ ($1 \mapsto 1/2$, $1/2 \mapsto 1/4$, $\dots$), and leave every other real number $x \notin S$ exactly where it is!
\end{enumerate}

\newpage
